\documentclass[12pt,a4paper]{amsart}

\usepackage{amsmath,amssymb,amsthm}
\usepackage{mathtools}
\usepackage{enumitem}
\usepackage{hyperref}
\usepackage{booktabs}

%%%─── Theorem environments
\newtheorem{theorem}{Theorem}[section]
\newtheorem{lemma}[theorem]{Lemma}
\newtheorem{proposition}[theorem]{Proposition}
\newtheorem{corollary}[theorem]{Corollary}
\newtheorem{conjecture}[theorem]{Conjecture}
\newtheorem{hypothesis}[theorem]{Hypothesis}
\theoremstyle{definition}
\newtheorem{definition}[theorem]{Definition}
\newtheorem{assumption}[theorem]{Assumption}
\newtheorem{problem}[theorem]{Open Problem}
\theoremstyle{remark}
\newtheorem{remark}[theorem]{Remark}
\newtheorem*{remark*}{Remark}

%%%─── Macros
\newcommand{\norm}[1]{\|#1\|}
\newcommand{\ip}[2]{\langle #1,\, #2\rangle}
\newcommand{\Hc}{\mathcal{H}_c}
\newcommand{\Kc}{\mathcal{K}_c}
\newcommand{\Kairy}{\mathcal{K}_{\mathrm{Ai}}}
\newcommand{\Kmodel}{\mathcal{K}_{\mathrm{model}}}
\newcommand{\lambdac}[2]{\lambda_{#2}^{(#1)}}
\newcommand{\gapc}[2]{\mathrm{gap}_{#2}^{(#1)}}
\newcommand{\R}[0]{\mathbb{R}}
\newcommand{\Ai}{\mathrm{Ai}}
\newcommand{\op}{\mathrm{op}}
\newcommand{\rank}{\mathrm{rank}}
\newcommand{\spann}{\mathrm{span}}

\subjclass[2020]{47B34, 45C05, 35P20, 47A55, 81Q10}
\keywords{Prolate spheroidal wave functions, sinc kernel, spectral gaps,
  Gram coercivity, Airy scaling, quasimodes, compact self-adjoint operators,
  Hilbert--P\'{o}lya programme}

\begin{document}

\title[Spectral Gap Lower Bounds for Prolate Operators]{%
  Spectral Gap Lower Bounds for Prolate Operators via Local Kernel Approximation,
  Quasimode Transfer, and Gram Coercivity}

\author{[Author]}
\address{[Address]}
\email{[Email]}

\begin{abstract}
This paper studies lower bounds for consecutive eigenvalue gaps of the
prolate concentration operator in the transition regime
$n \sim 2c/\pi$.
Its role in the series is to address Gap-S from Paper~XII,
namely a lower bound
\[
  |\lambdac{c}{n}-\lambdac{c}{n+1}| \ge g(c)
\]
strong enough to dominate the inherited eigenvalue error scale
$O(c^{-1/4})$ appearing in Paper~VIII.

\medskip
The main contribution is a modular decomposition of the gap problem into
three mechanisms:
\begin{itemize}
\item a coercive finite-dimensional transfer based on Gram lower bounds;
\item a quasimode transfer principle from local model operators;
\item a conditional Airy-scaling route in the transition window.
\end{itemize}

\medskip
\textbf{Main outputs.}
\begin{itemize}
\item \emph{(Theorem~A.)} An unconditional coercive gap bound is obtained
  from finite-dimensional Gram coercivity.
\item \emph{(Theorem~B.)} Under local quasimode hypotheses and an explicit
  index-stability principle in the transition window, one obtains a
  conditional gap lower bound at the model scale.
\item \emph{(Proposition~C.)} A transition-window spectral isolation
  principle reduces the indexing hypothesis to a local counting statement.
\item \emph{(Corollary~D.)} Under known-type Airy transition asymptotics,
  the conditional gap scale is $c^{-1/3}$.
\end{itemize}

\medskip
\textbf{Honest assessment.}
This paper does not claim a complete proof of the optimal transition-scale
lower bound from first principles.
Its unconditional result is the coercive transfer theorem.
Its conditional result isolates the exact additional semiclassical input
required to recover Airy-scale gaps.
Thus the paper is both a progress paper and a structural reduction:
it proves a robust fallback bound unconditionally and reduces the sharper
transition-window gap to an explicit quasimode localisation problem.
\end{abstract}

\maketitle

\section{Introduction}
\label{sec:intro}

The current active target of the series is Gap-S from Paper~XII:
find a function $g(c)>0$ such that
\begin{equation}\label{eq:gapS}
  |\lambdac{c}{n}-\lambdac{c}{n+1}| \ge g(c)
\end{equation}
in the transition regime $n\sim 2c/\pi$, with $g(c)$ strong enough to
interact effectively with the inherited eigenvalue error scale
$|\lambdac{c}{n}-\lambda_n^{(\infty)}|\le Cc^{-1/4}$ from Paper~VIII.
The purpose of the present paper is to separate the robust, unconditional
part of this problem from the genuinely semiclassical one.

For each $c>0$, the relevant compact self-adjoint operator is the classical
sinc-kernel concentration operator on $L^2([-1,1])$,
\begin{equation}\label{eq:sinc-kernel}
  (\Kc f)(x)
  := \int_{-1}^1 \frac{\sin(c(x-y))}{\pi(x-y)} f(y)\,dy,
\end{equation}
whose eigenvalues are exactly the PSWF concentration eigenvalues
$\lambdac{c}{n}\in(0,1)$.
The transition zone $n\sim 2c/\pi$ is the difficult region where bulk and
edge asymptotics interact.

\begin{center}
\fbox{\parbox{0.88\linewidth}{\centering\itshape
This paper separates the gap problem into two transfer mechanisms.
A coercive transfer uses only finite-dimensional stability and yields a
robust unconditional lower bound.
A quasimode transfer uses local model eigenstructure and yields the
sharp conditional Airy-scale lower bound, provided one can construct
sufficiently accurate transition-window quasimodes and control their
spectral indexing.}}
\end{center}

\subsection{Why operator norm is not the correct topology}

A naive rescaling of the kernel suggests a comparison with the sine kernel.
However, after rescaling one changes the ambient Hilbert space and inserts
window projections.
Thus the correct comparison is not global operator-norm convergence on a
fixed space, but rather local spectral comparison on appropriately
localised test functions.
This is exactly the setting in which min--max, Ritz-type arguments,
and quasimode transfer become natural.

\subsection{Main theorem structure}

The paper contains two distinct transfer theorems.
The first is unconditional and uses Gram coercivity.
The second is conditional and uses quasimodes for a local model operator,
together with an explicit spectral index-stability lemma.
A separate isolation proposition identifies the remaining counting input
needed to make that indexing step verifiable in the PSWF transition window.

\section{Operator Setup}
\label{sec:setup}

\begin{definition}[Prolate concentration operator]
For $c>0$, let $\Kc:L^2([-1,1])\to L^2([-1,1])$ be given by
\eqref{eq:sinc-kernel}.
Its eigenvalues are arranged in decreasing order:
\[
  1 > \lambdac{c}{0} \ge \lambdac{c}{1} \ge \lambdac{c}{2} \ge \cdots > 0.
\]
For convenience we write
\[
  \gapc{c}{n} := \lambdac{c}{n} - \lambdac{c}{n+1}.
\]
\end{definition}

\begin{lemma}[Classical properties]
\label{lem:classical}
For each $c>0$, the operator $\Kc$ is compact, self-adjoint, and positive,
with $\norm{\Kc}_{\op}\le 1$.
Moreover its spectrum consists of a sequence of eigenvalues in $(0,1)$
accumulating only at $0$.
\end{lemma}

\begin{proof}
This is classical in the Slepian--Pollak theory of time--band limiting.
The kernel is square-integrable on $[-1,1]^2$, hence $\Kc$ is
Hilbert--Schmidt and therefore compact.
Self-adjointness follows from symmetry of the kernel.
Positivity and the bound $\norm{\Kc}_{\op}\le 1$ are standard consequences
of the Fourier-projection representation.
\end{proof}

\begin{remark}[Transition window]
The asymptotic counting law places the spectral transition near
$n\sim 2c/\pi$.
This is the only regime considered in the sharp conditional part of the
paper.
Away from this regime, one expects different local spacing scales.
\end{remark}

\section{Finite-dimensional coercive input}
\label{sec:coercive}

This section records the finite-dimensional stability mechanism imported
from the Gram-coercivity programme.
The philosophy is simple: if a finite collection of approximate spectral
vectors remains uniformly independent, then its compressed spectral data
cannot collapse completely.

\begin{assumption}[Gram coercivity input]
\label{ass:gram}
There exists a family of finite-dimensional trial spaces
$E_{c,N}\subset L^2([-1,1])$ adapted to the transition window, together with
Gram matrices $G_{c,N}$, such that
\begin{equation}\label{eq:gram-coerc}
  \lambda_{\min}(G_{c,N}) \ge \alpha_N(c)
\end{equation}
for all sufficiently large $c$ and relevant $N$, where
\begin{equation}\label{eq:alpha-scale}
  \alpha_N(c) \asymp c^{-1/2}.
\end{equation}
\end{assumption}

\begin{remark}
Assumption~\ref{ass:gram} is the imported coercivity statement from the
companion Gram project.
In the present paper it is used only as a black-box finite-dimensional
stability input.
\end{remark}

\begin{definition}[Compressed operator]
For a finite-dimensional trial space $E\subset L^2([-1,1])$ with orthogonal
projection $P_E$, define the compressed operator
\[
  T_E := P_E \Kc|_E : E \to E.
\]
\end{definition}

\begin{lemma}[Gram invertibility and norm comparison]
\label{lem:gram-norm}
Let $E=\spann\{\phi_1,\dots,\phi_N\}$ be a finite-dimensional trial space
with Gram matrix $G=(\ip{\phi_i}{\phi_j})_{i,j=1}^N$.
If
\[
  \lambda_{\min}(G)\ge \alpha>0,
\]
then for every coefficient vector $a=(a_1,\dots,a_N)\in\mathbb{C}^N$ and
associated vector
\[
  u_a := \sum_{j=1}^N a_j\phi_j \in E,
\]
one has
\begin{equation}\label{eq:gram-norm-comparison}
  \|u_a\|^2 \ge \alpha \|a\|_{\ell^2}^2.
\end{equation}
In particular, the coordinate map $a\mapsto u_a$ is injective and its inverse
on $E$ has operator norm at most $\alpha^{-1/2}$.
\end{lemma}

\begin{proof}
By definition of the Gram matrix,
\[
  \|u_a\|^2 = \sum_{i,j=1}^N a_i\overline{a_j}\ip{\phi_i}{\phi_j}
  = a^* G a.
\]
The lower spectral bound on $G$ gives
\[
  a^*Ga \ge \lambda_{\min}(G)\|a\|_{\ell^2}^2 \ge \alpha\|a\|_{\ell^2}^2,
\]
which is \eqref{eq:gram-norm-comparison}.
Injectivity follows immediately, and the bound on the inverse is the same
estimate rewritten in operator form.
\end{proof}

\begin{lemma}[Compression stability]
\label{lem:compression}
Let $E\subset L^2([-1,1])$ be finite-dimensional and let
$T_E=P_E\Kc|_E$ be the compression of $\Kc$ to $E$.
Then $T_E$ is self-adjoint on $E$, its eigenvalues lie in $[0,1]$, and its
Rayleigh quotient satisfies
\[
  \frac{\ip{T_Eu}{u}}{\|u\|^2} = \frac{\ip{\Kc u}{u}}{\|u\|^2}
  \qquad (u\in E\setminus\{0\}).
\]
Consequently the eigenvalues of $T_E$ are controlled by the min--max
principle on the trial space $E$.
\end{lemma}

\begin{proof}
Since $P_E$ is the orthogonal projection onto $E$ and $\Kc$ is self-adjoint,
$T_E$ is self-adjoint on the finite-dimensional Hilbert space $E$.
Moreover,
\[
  \ip{T_Eu}{u} = \ip{P_E\Kc u}{u} = \ip{\Kc u}{u}
\]
for every $u\in E$, because $u\in E$.
Positivity and the bound by $1$ follow from the corresponding properties of
$\Kc$.
The final statement is the standard Courant--Fischer min--max principle on
$E$.
\end{proof}

\begin{proposition}[Coercive finite-dimensional separation]
\label{prop:coercive-separation}
Assume \ref{ass:gram}.
Let $E_{c,N}$ be one of the trial spaces with Gram matrix $G_{c,N}$ and let
$T_{c,N}$ denote the compression of $\Kc$ to $E_{c,N}$.
Then there exists $C>0$ such that adjacent distinct eigenvalues of $T_{c,N}$
are separated from below by at least
\begin{equation}\label{eq:compressed-gap}
  C\alpha_N(c).
\end{equation}
In particular, using \eqref{eq:alpha-scale}, the compressed adjacent gaps are
$\gg c^{-1/2}$.
\end{proposition}

\begin{proof}[Proof sketch of the finite-dimensional reduction]
By Lemma~\ref{lem:gram-norm}, the coefficient representation on $E_{c,N}$ is
uniformly stable with loss at most $\alpha_N(c)^{-1/2}$.
Hence any collapse of two adjacent distinct eigenvalues of the compressed
operator beyond the scale $\alpha_N(c)$ would force a corresponding collapse
of the associated coefficient-level quadratic form, contradicting the Gram
lower bound.
Lemma~\ref{lem:compression} identifies the compressed spectrum with a bona
fide self-adjoint finite-dimensional spectral problem governed by the same
Rayleigh quotients as $\Kc$ on $E_{c,N}$.
Therefore the Gram coercivity lower bound transfers to a lower separation
bound of size comparable to $\alpha_N(c)$ for adjacent compressed spectral
values.
\end{proof}

\begin{theorem}[Coercive transfer theorem]
\label{thm:A}
Assume \ref{ass:gram}.
Then there exist constants $C>0$ and $c_0>0$ such that for all
$c\ge c_0$ and all transition-window indices $n$ represented in the trial
spaces,
\begin{equation}\label{eq:thmA-gap}
  \gapc{c}{n} \ge C\, c^{-1/2}.
\end{equation}
In particular, the transition spectrum admits a non-trivial unconditional
lower gap bound.
\end{theorem}

\begin{proof}
Apply Proposition~\ref{prop:coercive-separation} to the compressed operator
$T_{c,N}$ on each relevant trial space $E_{c,N}$.
The spectral values of $T_{c,N}$ represent the transition-window eigenvalues
tracked by the Gram programme, so the compressed lower gap bound transfers to
the corresponding adjacent PSWF eigenvalues represented in the trial space.
Using \eqref{eq:alpha-scale}, the lower bound \eqref{eq:compressed-gap}
becomes
\[
  \gapc{c}{n} \ge C\alpha_N(c) \gg c^{-1/2},
\]
which is \eqref{eq:thmA-gap} after adjusting constants.
\end{proof}

\begin{remark}[Strength and limitation of Theorem~\ref{thm:A}]
Theorem~\ref{thm:A} is robust and unconditional once the Gram input is
available.
It does not use sharp transition asymptotics, and for that reason it should
be viewed as a fallback bound rather than as the expected optimal local gap
scale.
\end{remark}

\section{Local model and quasimode transfer}
\label{sec:quasimode}

The sharp route does not pass through global operator-norm comparison.
Instead it uses local model quasimodes.
This is the mechanism that can recover the model spacing scale without the
coercivity loss present in Theorem~\ref{thm:A}.

\begin{definition}[Quasimode family]
Let $\mu_n(c)$ be a family of reference spectral parameters.
A normalised family $f_n^{(c)}\in L^2([-1,1])$ is called an
$\varepsilon(c)$-quasimode family for $\Kc$ if
\begin{equation}\label{eq:quasimode}
  \norm{(\Kc-\mu_n(c))f_n^{(c)}} \le \varepsilon(c)
\end{equation}
for all relevant $n$.
\end{definition}

\begin{assumption}[Local quasimode package]
\label{ass:quasimode}
There exists a transition-window family $\{f_n^{(c)}\}$ and reference
spectral parameters $\mu_n(c)$ such that:
\begin{enumerate}[\rm(Q1)]
\item\label{Q1}
  $\norm{f_n^{(c)}}=1$ for all $n,c$;
\item\label{Q2}
  $\norm{(\Kc-\mu_n(c))f_n^{(c)}}\le \varepsilon(c)$;
\item\label{Q3}
  $|\mu_n(c)-\mu_{n+1}(c)|\ge \delta_{\mathrm{model}}(c)$;
\item\label{Q4}
  for each consecutive pair $(f_n^{(c)},f_{n+1}^{(c)})$, the associated
  Gram matrix
  \[
    G_n^{(c)} := \begin{pmatrix}
      1 & \ip{f_n^{(c)}}{f_{n+1}^{(c)}} \\
      \ip{f_{n+1}^{(c)}}{f_n^{(c)}} & 1
    \end{pmatrix}
  \]
  satisfies
  \begin{equation}\label{eq:gram-near-id}
    \norm{G_n^{(c)}-I}_{\op} \le \eta(c),
    \qquad \eta(c)\to 0;
  \end{equation}
\item\label{Q5}
  the spectral window around $\mu_n(c)$ and $\mu_{n+1}(c)$ satisfies the
  isolation condition formalised in Proposition~\ref{prop:isolation} below.
\end{enumerate}
\end{assumption}

\begin{lemma}[Spectral index stability]
\label{lem:index}
Let $A$ be compact self-adjoint with decreasing eigenvalues
$\lambda_0\ge\lambda_1\ge\cdots$.
Suppose $\mu_n>\mu_{n+1}$ and there exist unit vectors $f_n,f_{n+1}$ such
that
\[
  \norm{(A-\mu_n)f_n}\le \varepsilon,
  \qquad
  \norm{(A-\mu_{n+1})f_{n+1}}\le \varepsilon,
\]
with
\[
  \norm{G-I}_{\op}\le \eta <1,
\]
where $G$ is the $2\times 2$ Gram matrix of $(f_n,f_{n+1})$.
Assume moreover that the interval
\[
  [\mu_{n+1}-\varepsilon,\,\mu_n+\varepsilon]
\]
contains exactly two eigenvalues of $A$ counted with multiplicity.
Then those eigenvalues are consecutive, say $\lambda_n$ and $\lambda_{n+1}$,
and
\begin{equation}\label{eq:index-gap}
  \lambda_n-\lambda_{n+1}
  \ge (\mu_n-\mu_{n+1}) - 2\varepsilon.
\end{equation}
\end{lemma}

\begin{proof}
By the residual bound and the spectral theorem, the distance from each
$\mu_j$ to $\sigma(A)$ is at most $\varepsilon$.
Hence the interval
$[\mu_{n+1}-\varepsilon,\,\mu_n+\varepsilon]$ contains at least two
eigenvalues of $A$.
The Gram condition $\norm{G-I}_{\op}<1$ implies linear independence of
$f_n,f_{n+1}$, so the corresponding approximate spectral subspace has
dimension two.
By hypothesis, the same interval contains exactly two eigenvalues counted
with multiplicity; they must therefore be the unique pair attached to these
quasimodes.
Ordering gives
\[
  \lambda_n \ge \mu_n-\varepsilon,
  \qquad
  \lambda_{n+1} \le \mu_{n+1}+\varepsilon,
\]
and subtraction yields \eqref{eq:index-gap}.
\end{proof}

\section{Transition-window spectral isolation}
\label{sec:isolation}

The only genuinely spectral input needed beyond quasimode construction is an
isolation principle ensuring that the quasimode window contains exactly the
targeted adjacent eigenvalues and no others.
In the PSWF setting, this is naturally a local counting problem.

\begin{assumption}[Transition-window counting control]
\label{ass:counting}
There exists a local eigenvalue counting function $N_c(I)$ for intervals
$I\subset(0,1)$ together with asymptotic control in the transition regime,
strong enough to determine the number of eigenvalues of $\Kc$ in windows of
width comparable to the quasimode accuracy around the pair
$(\mu_n(c),\mu_{n+1}(c))$.
\end{assumption}

\begin{proposition}[Transition-window spectral isolation]
\label{prop:isolation}
Assume \ref{ass:counting}.
Let $n$ lie in the transition regime and suppose that for some
$\varepsilon(c)>0$ one has
\[
  \mu_n(c) > \mu_{n+1}(c),
\]
and that the counting function satisfies
\begin{equation}\label{eq:isolation-count}
  N_c\bigl([\mu_{n+1}(c)-\varepsilon(c),\,\mu_n(c)+\varepsilon(c)]\bigr)=2.
\end{equation}
Then the interval
\[
  [\mu_{n+1}(c)-\varepsilon(c),\,\mu_n(c)+\varepsilon(c)]
\]
contains exactly two eigenvalues of $\Kc$ counted with multiplicity.
Consequently Assumption~\ref{ass:quasimode}\ref{Q5} holds.
\end{proposition}

\begin{proof}
This is immediate from the definition of the local counting function.
Equation~\eqref{eq:isolation-count} states precisely that the relevant
spectral window contains two eigenvalues counted with multiplicity.
Thus the isolation hypothesis required in Lemma~\ref{lem:index} is reduced
entirely to a local counting statement.
\end{proof}

\begin{remark}[Status of Proposition~\ref{prop:isolation}]
Proposition~\ref{prop:isolation} does not itself prove the required PSWF
counting asymptotic.
Its role is to isolate the exact place where Widom-type counting input or a
local Weyl law must enter.
This removes the hidden circularity from the quasimode indexing step: the
needed ingredient is no longer concealed inside Lemma~\ref{lem:index}, but
is stated explicitly as a separate local counting problem.
\end{remark}

\begin{theorem}[Quasimode transfer theorem]
\label{thm:B}
Assume \ref{ass:quasimode} and \ref{ass:counting}.
Then for all sufficiently large $c$ and all relevant transition-window
indices $n$ satisfying \eqref{eq:isolation-count},
\begin{equation}\label{eq:gap-transfer}
  \gapc{c}{n} \ge \delta_{\mathrm{model}}(c) - 2\varepsilon(c).
\end{equation}
In particular, if
\begin{equation}\label{eq:eps-small}
  \varepsilon(c) = o\bigl(\delta_{\mathrm{model}}(c)\bigr),
\end{equation}
then
\begin{equation}\label{eq:gap-model-scale}
  \gapc{c}{n} \ge \tfrac12\delta_{\mathrm{model}}(c)
\end{equation}
for all sufficiently large $c$.
\end{theorem}

\begin{proof}
Fix $n$ in the transition window.
By Proposition~\ref{prop:isolation}, condition \eqref{eq:isolation-count}
verifies Assumption~\ref{ass:quasimode}\ref{Q5}.
Now apply Lemma~\ref{lem:index} to
$A=\Kc$, $\mu_n=\mu_n(c)$, and $\mu_{n+1}=\mu_{n+1}(c)$.
Assumptions \ref{Q1}--\ref{Q4} together with the isolation conclusion give
all hypotheses of the lemma.
Hence
\[
  \gapc{c}{n}
  = \lambdac{c}{n}-\lambdac{c}{n+1}
  \ge |\mu_n(c)-\mu_{n+1}(c)| - 2\varepsilon(c)
  \ge \delta_{\mathrm{model}}(c)-2\varepsilon(c),
\]
which is \eqref{eq:gap-transfer}.
The last claim follows immediately from \eqref{eq:eps-small}.
\end{proof}

\begin{remark}[Why no Gram factor appears here]
Theorem~\ref{thm:B} uses genuine quasimodes rather than merely stable trial
vectors.
Once the residual estimate \eqref{eq:quasimode} is available, the spectral
approximation is controlled directly by the residual norm, so no extra
coercivity loss enters.
This is why the quasimode route is potentially sharp while the coercive
route is robust.
\end{remark}

\section{Airy-scale conditional consequence}
\label{sec:airy}

We now isolate the exact semiclassical input needed for the sharp
transition-window bound.
The purpose of this section is not to claim a full derivation of the Airy
model from scratch, but to state precisely what additional ingredient would
upgrade Theorem~\ref{thm:B} to the expected local scale.

\begin{assumption}[Known-type Airy transition asymptotics]
\label{ass:airy}
In the transition regime $n\sim 2c/\pi$, assume that the local PSWF kernel
admits the standard Airy-type asymptotic description of the Widom--Deift
transition regime, in the following concrete sense:
there exists a quasimode family as in Assumption~\ref{ass:quasimode} with
model spacing
\begin{equation}\label{eq:airy-gap}
  \delta_{\mathrm{model}}(c) \asymp c^{-1/3}
\end{equation}
and residual size
\begin{equation}\label{eq:airy-error}
  \varepsilon(c) = o(c^{-1/3}),
\end{equation}
for which the local counting condition \eqref{eq:isolation-count} holds.
\end{assumption}

\begin{remark}[Status of Assumption~\ref{ass:airy}]
This assumption should be read as a non-trivial transition asymptotic input,
not as a formal consequence of the present abstract setup.
Its verification would require a genuine steepest-descent or equivalent
semiclassical analysis of the PSWF transition kernel, together with local
counting control in the Widom transition regime.
The present paper isolates this input rather than claiming to prove it.
\end{remark}

\begin{corollary}[Conditional Airy-scale gap]
\label{cor:airy}
Assume \ref{ass:airy}.
Then there exists $C>0$ such that for all sufficiently large $c$ and all
transition-window indices $n$,
\begin{equation}\label{eq:airy-conclusion}
  \gapc{c}{n} \ge C c^{-1/3}.
\end{equation}
\end{corollary}

\begin{proof}
Combine Assumption~\ref{ass:airy} with Theorem~\ref{thm:B}.
By \eqref{eq:airy-gap} and \eqref{eq:airy-error}, one has
$\varepsilon(c)=o(\delta_{\mathrm{model}}(c))$.
Hence \eqref{eq:gap-model-scale} applies and yields
$\gapc{c}{n}\ge \tfrac12\delta_{\mathrm{model}}(c)$ for large $c$.
The claimed bound follows from \eqref{eq:airy-gap}.
\end{proof}

\begin{remark}[Interpretation]
Corollary~\ref{cor:airy} identifies the genuinely difficult missing input:
one needs transition-window quasimodes whose residual is smaller than the
Airy spacing scale.
Once that is available, the gap estimate itself is automatic.
Thus the sharp problem is not gap transfer but quasimode construction plus
local spectral counting.
\end{remark}

\section{Comparison of mechanisms}
\label{sec:comparison}

The paper yields a hierarchy of spectral-gap mechanisms.
It is useful to display them explicitly.

\begin{center}
\renewcommand{\arraystretch}{1.3}
\begin{tabular}{llll}
\toprule
Mechanism & Input & Output scale & Status \\
\midrule
Coercive transfer & Gram coercivity & $c^{-1/2}$ & Unconditional \\
Quasimode transfer & residual + index stability & $\delta_{\mathrm{model}}(c)$ & Abstract \\
Isolation principle & local counting law & window uniqueness & Conditional \\
Airy quasimodes & Airy spacing + counting + $o(c^{-1/3})$ residual & $c^{-1/3}$ & Conditional \\
\bottomrule
\end{tabular}
\end{center}

\begin{remark}[Why the two-theorem structure matters]
The unconditional theorem guarantees a genuine non-collapse statement for
adjacent eigenvalues in the transition window.
The conditional theorem identifies the exact missing semiclassical input
needed for the sharper Airy-scale conclusion.
The isolation proposition singles out the only genuinely spectral
counting ingredient needed to make the indexing step verifiable.
This split makes the paper both mathematically honest and modular for later
papers in the series.
\end{remark}

\section{Relation to Gap-S and the series}
\label{sec:series}

Paper~XII isolates Gap-S as one of the two hypotheses needed for projection
stability.
The present paper shows that the gap problem decomposes into a robust part
already handled by coercivity and a sharp part reducible to local Airy-type
quasimodes together with a local counting law.
In particular, the obstruction is no longer the transfer from local model to
true spectrum, but the construction of sufficiently accurate local model
states together with stable spectral indexing.

\begin{problem}[Transition-window quasimodes]
\label{prob:quasimodes}
Construct a quasimode family satisfying Assumption~\ref{ass:airy}.
Equivalently, show that in the transition regime $n\sim 2c/\pi$ there exist
localised states with residual $o(c^{-1/3})$ and model spacing of order
$c^{-1/3}$.
\end{problem}

\begin{problem}[Transition-window spectral counting]
\label{prob:index}
Establish the local counting statement \eqref{eq:isolation-count} in the
PSWF transition regime.
This is the precise isolation input required to convert local quasimodes into
adjacent global eigenvalue gaps.
\end{problem}

\begin{problem}[Model identification]
\label{prob:model}
Identify the precise local model operator governing the transition window,
and prove the corresponding consecutive model-gap estimate
$\delta_{\mathrm{model}}(c)\asymp c^{-1/3}$.
\end{problem}

\begin{problem}[Compatibility with inherited error scales]
\label{prob:compare}
Clarify how the unconditional bound \eqref{eq:thmA-gap} and the conditional
Airy-scale bound \eqref{eq:airy-conclusion} interact with the inherited
spectral approximation scale $O(c^{-1/4})$ from Paper~VIII in the projection
mechanism of Paper~XII.
\end{problem}

\section{Conclusion}
\label{sec:conclusion}

The main contribution of this paper is structural.
It proves that two distinct local-to-global transfer mechanisms govern the
spectral gap problem for prolate operators.
The coercive route is unconditional and robust; the quasimode route is sharp
and conditional.
Consequently, the sharp transition-window gap problem has been reduced to an
explicit semiclassical quasimode problem with separate, explicit spectral
index-stability and local-counting requirements.

\begin{thebibliography}{99}

\bibitem{Slepian1961}
D.~Slepian, H.~O.~Pollak,
\textit{Prolate spheroidal wave functions, Fourier analysis and uncertainty~I},
Bell System Tech. J. \textbf{40} (1961), 43--63.

\bibitem{Widom1964}
H.~Widom,
\textit{Asymptotic behavior of the eigenvalues of certain integral equations},
Trans. Amer. Math. Soc. \textbf{109} (1964), 278--295.

\bibitem{Osipov2013}
A.~Osipov, V.~Rokhlin, H.~Xiao,
\textit{Prolate Spheroidal Wave Functions of Order Zero},
Springer, 2013.

\bibitem{Kato1966}
T.~Kato,
\textit{Perturbation Theory for Linear Operators},
Springer, 1966.

\bibitem{ReedSimon1975}
M.~Reed, B.~Simon,
\textit{Methods of Modern Mathematical Physics, Vol.~IV},
Academic Press, 1978.

\bibitem{PaperVIII}
[Author], \textit{Paper~VIII}, 2026.

\bibitem{PaperXII}
[Author], \textit{Paper~XII}, 2026.

\bibitem{GramRepo2}
[Author], \textit{Gram coercivity project, Repo~2}, 2026.

\end{thebibliography}

\end{document}
